\section{CASSCF-Type Wave Functions}\label{sec-diabat-wfn-casscf}
A CASSCF state is a linear combination of determinants within a specified active space,
\begin{equation}
|\Psi_I\rangle=\sum_{D\in\mathrm{CAS}}C_D^{(I)}|D\rangle.
\end{equation}
Orbitals outside the active space have fixed occupations. Accordingly, the state file needs the orbital-space definition and, for each state, the active-orbital occupations and their CI coefficients.

\subsection*{Complete CASSCF example}
The following example has two active orbitals and two active electrons. Both determinants have one alpha and one beta electron. The numerical values are illustrative.
\begin{lstlisting}[style=technicalfile]
Wave-function data exported by Diabat
Method: CASSCF
Number of active orbitals: 2
Number of active electrons: 2
Number of doubly occupied orbitals: 0
Number of doubly virtual orbitals: 0
Number of frozen occupied orbitals: 0
Number of frozen virtual orbitals: 0
Number of electrons: 2
ID of the first doubly occupied orbital: 1
ID of the first active orbital: 1
ID of the first doubly virtual orbital: 3

State Information:
State index: 1
Total energy (a.u.): -1.0
Excitation energy (eV): 0.0
<S^2>: 0.0
  Amplitudes:
 ---------------------------------
   [20]    7.071067811865475E-01
   [02]    7.071067811865475E-01
 ---------------------------------
\end{lstlisting}
The strings \texttt{20} and \texttt{02} describe two different doubly occupied active orbitals. Their squared coefficients sum to one. The associated MO file must contain the two active orbitals in the order specified by the header.

\subsection*{Occupation strings}
A CASSCF coefficient record consists of a bracketed occupation string followed by one real coefficient. Each character describes one active \emph{spatial} orbital, in order:
\begin{center}
\begin{tabular}{cl}
\hline
Character & Occupation\\\hline
\texttt{2} & One alpha and one beta electron\\
\texttt{u} & One alpha electron\\
\texttt{d} & One beta electron\\
\texttt{0} & Empty orbital\\\hline
\end{tabular}
\end{center}
For instance, \texttt{[ud]} denotes an alpha electron in the first active orbital and a beta electron in the second. Each string must have exactly as many characters as there are active orbitals. All determinants belonging to a state must have the same numbers of alpha and beta electrons. Unlisted configurations have zero coefficients.

The doubly occupied and doubly virtual orbital counts describe the blocks outside the active space. The corresponding \texttt{ID of the first ... orbital} fields give their starting MO indices. Frozen-orbital counts are included as separate header fields. All counts and indices must agree with the MO file.

\subsection*{Coefficient signs}
The CASSCF representation uses the spin-orbital ordering \(\alpha_1,\beta_1,\alpha_2,\beta_2,\ldots\). If the source program instead orders all alpha orbitals before all beta orbitals, convert its determinant coefficients to this interleaved ordering. For alpha and beta occupation numbers \(n_{\alpha i}\) and \(n_{\beta i}\), the required sign is
\begin{equation}
(-1)^{\sum_i n_{\beta i}\sum_{j>i}n_{\alpha j}}.
\end{equation}
For example, \texttt{[du]} changes sign between the two ordering conventions, whereas \texttt{[ud]} does not.
